% 原书第 131 页（PDF 第 144 页）
\BookPage{131}{144}
\BookHeading{十七世纪民间天道观研究—基于《广东新语》的研读}{十七世纪民间天道观研究\\[3pt]—基于《广东新语》的研读}{周正庆*}
「天道观」指人们对于世界本原的观点，广义包括人们围绕着「天」展开对天人关系、自然现象的系列解读而形成的观点。

中国古代的天道观最早是一个哲学命题，殷周之际的《易经》将八种自然现象：天、地、雷、风、水、火、山、泽，分别设计成干、坤、震、巽、坎、离、艮、兑八种方位，以金木水火土五行相生相克理论进行解读，说明宇宙的生成及万物间的联系和变易，形成了系统化的阴阳八卦，五行相生相克天道观，并对后世造成了巨大的影响。

在论述十七世纪中国民间天道观之前，先让我们看一下同时代欧洲社会的天道观。文艺复兴前欧洲封建社会的天道观长期被宗教神学所垄断，而封建宗教的天道观赖以创立的理论是托勒密创立于西元二世纪的「地心说」，这种封建桎梏下的天道观直到十四世纪欧洲文艺复兴开始才受到冲出。

十五世纪的波兰天文学家哥白尼在《天体运行论》一书中，用科学方法描述了地球的运动和宇宙的构造情形，推翻了影响欧洲一千多年的「地心说」，发展成为较为科学的「日心说」，也从根本上动摇了欧洲唯心主义宇宙观的理论基础，至西元1600年2月17日，经过八年监禁后的义大利哲学家布鲁诺由于支持哥白尼的理论而在罗马百花广场惨遭火刑，但欧洲的科学家并没有因此而停止前进的脚步，1609年，伽俐略创制了天文望远镜并用来观测天体，找到了确凿支持哥白尼的学说证据，欧洲人的天文观得以挣脱了长期占统治地位神学的桎梏，至十七世纪中叶，牛顿以数量推理和实证科学为主体取得了一系列科学成就，从此十七世纪西方天文

\vfill
\begin{BookNotes}
\BookNote{*}{作者任职于广州暨南大学历史系。}
\end{BookNotes}

% 原书第 132 页（PDF 第 145 页）
\BookPage{132}{145}
\noindent 学进入了现代科学体系的发展轨道，由此产生了较为科学的天道观，知道天体运行的一般规律，并对风雨雷电等自然现象有了较为正确对待的认识，科学的发展带来民间天道观科学的重建在欧洲先进国家得以普及，十七世纪到中国的传教士所带来的关于天道的科学理论就已经说明了这一点。

相对于欧洲天文学的快速发展带来天道观的改变，中国人的天道观又是如何的呢？我们可以通过对《广东新语》的解读寻得答案。

屈大均，广东番禺人，生于明崇祯三年，卒于康熙三十五年（1630-1696），其活动的时间段基本上是十七世纪的时间范畴内，而屈大均是明朝晚年的诸生，也即是明末清初知识分子的代表，屈大均的一生并没有对历法作专门的研究，所以对天体的运行并没有专门的知识，故而笔者认为屈大均所着《广东新语》所反映出来的天道观代表着民间的天道观。

\begin{BookQuote}
《广东新语》是屈大均「考方舆，披志乘，验之以身经，征之以目睹」2以亲身
\end{BookQuote}

之经历为基础，于康熙二十六年（1687）写成的一本笔记体裁之作3，此书以语为纲，共为天、地、山、水、石、神、人、女、事、学、文、诗、艺、货、器、宫、舟、坟、禽、兽、鳞、介、虫、木、香、草、怪二十八语，一语一卷，共二十八卷，卷下设目，总计869目，主要记述了屈大均生活的十七世纪及以前广东的天地山川、风土人情、物产名胜、神话传说等内容，其中涉及到天道观内容集中在卷一至卷六的天、地、山、水、石、神等六语里，本文以这六卷作为重点对屈大均的天道观进行分析。

\vfill
\begin{BookNotes}
\BookNote{1.}{本文采用的版本为中华书局1985年版。}
\BookNote{2.}{（清）屈大均着，《广东新语·潘序》，中华书局，1978年。}
\BookNote{3.}{1关于《广东新语》一书的成书年代，学术界有争论，汪宗衍、赵立人、南炳文等人均有不同说法，笔者采用吴建新先生康熙二十六年之说，详见吴建新：<《广东新语》成书年期再探〉一文，载《广东社会科学》，1989年第3期。}
\end{BookNotes}

% 原书第 133 页（PDF 第 146 页）
\BookPage{133}{146}
\BookSubheading{一、从《广东新语》看屈大均的天道观}

\BookSubheading{（一）世界本原观}

关于天地的构成问题，屈氏在书中没有作专门的论述，但透过其对广东地理山川的记载及针对地理山川的形成而引发的解释，我们可以对屈大均的世界本原观点进行客观的分析。

屈大均的世界本原论主要体现在对地球的认识上，《广东新语》一书中多处反映了其「地圆天方」观。

在十七世纪的欧洲，地球作为一个球体这一基本知识已经为一般欧洲知识分子所掌握，但屈大均在《广东新语》中仍然用「地圆天方」的传统理论去解释自然现象，书中多次提到粤地是地之尽头，即是其持「地方论」的明证，如其在《广东新语·地语》中认为「地至广南而尽，尽者，尽之于海也」。海南岛只不过是「广南之余地在海中者也4，就是这种观点的反映。

又比如，屈大均在论及广东的飓风时，对广东民谚「南风尾，北风头」作了这样的解释，认为广东「盖地势北高而南下，南风从下而起，故为尾。又南风愈吹愈大，其大在尾。北风初起即大，其大在头故也」。从中我们可以看出屈大均认为冬天广东刮南风，大风在后头，而南为广东地势之尾，故其大在尾。而刮北风则相反，情形是「北风初起即大」，因为北方是广东地势之头，故「其大在头故也」。5对大地的构成作头尾的论述形象地反映了屈氏「地圆天方」的观点。

\BookSubheading{（二）自然现象观}

屈氏书中所论，对许多自然现象有正确的记述，但仅止于「经验主义」，没有正确的科学理论作指导，所以其解释最终流于唯心主义。

\begin{BookQuote}
飓风」是广东常见的自现现象，对此屈大均作了记载：南海岁有旧风，亦曰风旧，盖飓风也。其起也，自东北者必自北而西，自西北
\end{BookQuote}

\vfill
\begin{BookNotes}
\BookNote{4.}{同注1，卷2地语·地〉，第29页。}
\BookNote{5.}{同注1，卷1〈天语·风〉，第10页。}
\end{BookNotes}

% 原书第 134 页（PDF 第 147 页）
\BookPage{134}{147}
\begin{BookQuote}
者必自北而东，而俱至南乃息，谓之落西，亦曰荡西，又曰回南，凡二昼夜乃息，亦曰风痴。若不落西，不回南，则逾月复作，作必对时，日作次日止，夜作次夜止。谚曰：「朝北暮南子夜东。又曰：「朝三暮七，昼不过一。盖其暴者不久，或数时，或一日夜。其柔者久，或二三夜。有一岁再三作者，有三四岁不作者。6
\end{BookQuote}

屈氏的这段记载总结了飓风活动的几个规律，这些规律已经被实践验证是正确的：第一，认为台风的运行是不定向的，走的是「朝北暮南子夜东」的轨迹，但飓风的吹刮又是有规律的，无论是起于「东北者」或「西北者」最终以刮南风而结束，否则风仍未刮完，一月后必将再有飓风。第二，「飓风」再发也有一定的规律可循，那就是「日作次日止，夜作次夜止」。第三，关于「飓风」持续的时间，屈大均作了总结，认为猛烈的飓风「其暴者不久，或数时，或一日夜」，相反较和缓的飓风「其柔者久，或二三夜」。

屈大均这些解释尽管是切合实际的，但他在书中并没有作科学的解释，如前所述，屈氏将地「南风尾，北风头」现象归于「（地）势北高而南下」，将后者用民谚进行述说，这都不是科学的解释法。

对于一些自然现象，屈大均的解释虽然有科学的成分，但囿于知识结构所限没有作进一步科学述说，反而走向了唯心主义了。比如屈大均对广东地区海潮的形成及其特点进行了正确的叙述，认为广东的潮汐「日一潮而一汐，或日两潮而两汐」皆为正常，但有些像发生在番禺的潮水其情形看似不正常，但实际也非常常见：「朝潮未落，暮潮乘之，驾以终风，前后相蹙，海水为之沸溢，是曰遝潮。一岁有之，或再岁有之，此则潮之变，水之不能其节者也。若以岁之十月，自朔至于十有二日候潮，朔日潮盛，则明年正月必有大水。二日则应二月，日直其月，至于十有二日皆然」。对于这种现象，屈大均知道潮汐的产生与变化与月亮有关，月亮的圆缺和升降对潮汐产生很大的影响，其认为「月生明则潮初上，月中则潮平，月转则潮渐退，月没则潮干。月与日会，则潮随月而会，月与日对，则潮随月而对」。但当屈大均再次对月亮与潮水关系作进一步解释时又流于诸子的混元一气论和阴阳八

\vfill
\begin{BookNotes}
\BookNote{6.}{（清）屈大均着，《广东新语》，卷1〈天语·旧风〉，中华书局，1978年，第11页。}
\end{BookNotes}

% 原书第 135 页（PDF 第 148 页）
\BookPage{135}{148}
\noindent 卦论，未能对二者之间的关系作出合乎科学的解释，以至最后的解释没有取得合乎科学质的突破，其认为:「月者水之精，潮者月之气，精之所至，气亦至焉，此则水之常节也。盖水与月同一坎体，故以月为节者，在在有常，而以日为节者，在在有变也。余靖云：月之所临，则水往从之。故月临卯酉，则水涨乎东西。月临子午，则潮平乎南北，彼竭此盈，皆系于月不系于日是也。」

又比如对于涨潮的原因，屈大均还是将之归为「气」所为，最后不免流于唯心主义：

\begin{BookQuote}
涨者嘘吸先天之气以为升降，气升则长，长则潮下虚，下虚十丈。则潮上赢十丈，气降则消，消则潮下实，下实一尺，则潮上缩一尺，皆气之所为，故曰涨。凡水能实而不能虚，惟涨海虚时多而即时少，气之最盛故涨，若夫飓风发而咸流逆起，大伤禾稼，则气郁抑而不得其平，亦涨之说也。涨海故多飓风，故其潮信无定。雷州调黎之东，一日两潮而两汐，西则一。调黎之潮，东消而西长，那黄则西消而东长，其风不定，其潮汐因之，风者气之所鼓者也。平常则旧潮未去，新潮复来，常羡溢而不平，故曰涨海也。
\end{BookQuote}

\BookSubheading{（三）天人关系观}

屈大均在《广东新语》中，许多地方在记述自然现象和山川名胜时，体现出来的是一种天人感应唯心主义的思想，卷一〈天语·星〉便是这种思想的最明显体现，笔者选择了能代表屈氏观点的记载进行罗列，如下表所示：

卷一<广东新语·天语·星>中天人感应事件概况表8

\begin{BookQuote}
时间异象对应事件
\end{BookQuote}

秦二世二年五星会东井，倍于南斗秦失王气，则越得霸气，故尉佗应之而兴也梁大同五年冬彗出南斗，东南指，长一丈余李贲称帝交州

有二枉矢出北斗魁，委曲蛇形隋大业十二年秋九月豫章刘士弘自称楚帝，地及番禺

见于南斗

\vfill
\begin{BookNotes}
\BookNote{7.}{同注1，卷4〈水语·潮〉，第133页。}
\BookNote{8.}{（清）屈大均着，《广东新语》，卷1〈天语·星〉，中华书局，1978年，第5-7页。}
\end{BookNotes}

% 原书第 136 页（PDF 第 149 页）
\BookPage{136}{149}
\noindent 宋皇元年秋九月乙巳太白犯南斗广源蛮侬智高反祥兴元年秋八月乙巳星陨于海如雨师溃崖山，丞相陆秀夫抱帝赴海死永乐十三年夏有星孛于南斗南海盗黄萧养起，僭称东阳王，攻犯广州

屈大均认为「南斗固越之司命也」是广东的主星，主星受到侵犯或出现异象，必有不祥之事发生。如上表所显示，屈大均不惜用大量的篇章与事件说明他的这个观点。

我们再看〈天语·星聚〉，屈大均在论及明代广东文化兴盛时，认为斗牛与中星是广东文化之星，「斗牛与中星明，则其地儒道大兴」，斗牛星旺则广东文化必兴，为了说明「天象诚不虚垂示也」的观点，屈大均在书中罗列了许多星聚的现象，比如:洪武、永乐间，「五星两聚牛斗」，结果是「白沙先生出」。成化丙戌年间，「中星明于越之分野」，对应的是「甘泉以是岁生」、「自此粤士大夫多以理学兴起，肩摩砺接，彬彬乎有邹鲁之风」9。

关于世界的本原问题，民间主要反映在对天地构成的看法，有学者认为屈氏的宇宙观由以「心」学为主，淡而转向「朴素的唯物主义宇宙观」10，笔者认为如此评价有失偏颇。如上所述，笔者认为从《广东新语》所反映出屈氏天道观来看，仍然没有超过《易经》与先秦诸子的理论，体现出来的仍然是「地圆天方」的思想，对自然现象的看法仍以唯心者居多，并没作唯物之转变。

\BookSubheading{二、十七世纪民间天道观的评估}

对于十七世纪民间天道观的评估问题，我们必须看到此时适逢明清易代，而二个朝代的统治者对待民间对「天道」的探索探取了二种迥然不同的态度，导致民间天道观也产生了不同的结果，故笔者将之分为前期与后期进行探索，前者以明为主，后者则是入清之事了。

\vfill
\begin{BookNotes}
\BookNote{9.}{同注1，卷1〈天语·星聚〉，第8页。}
\BookNote{10.}{关汉华、冼剑民，〈屈大均及其史学〉，《暨南学报（哲学社会科学）》，1997年2期。}
\end{BookNotes}

% 原书第 137 页（PDF 第 150 页）
\BookPage{137}{150}
如上所述，屈大均的天道观思想代表着十七世纪前期即明代普通民众对「天」与「天」相关自然现象的看法，其观点基本停留在先秦的阶段，没有取得质的变化，但对于天体的认识，明代专业的科学家比民间百姓有着更为深刻与科学的认识，最大的体现在于历法的修订所反映出来天道观的进步。

明人开始接触到西方的天文知识，来源于义大利的耶稣会传教士利玛窦（MatteoRicci,1552-1610），利玛窦于明朝万历年间来到中国，着有《乾坤体义》一书，对托勒密（ClandiusPtolemaeus，约90-168）的「旧地心说」，亚理士多德（Aristotle,B.C.384-B.C.322）的四元行论」进行了介绍，是西方传教士第一次将欧洲天文学介绍到了中国，并引起了中国上层社会的极大兴趣。

由于受到西方传教士的影响，中国的科学家开始接触到西方的天文学，对于「天」的认识取得了一定的进步，徐光启对中国历法的修订与对天体的认识所取得的进步就是最明显的例子。

明末徐光启受命修订历法，将西方历法在《崇祯历书》中进行解读，负责对全书进行「释义演文，讲究润色，校勘试验」的编译。《崇祯历书》的编译，自崇祯四年（1631）起直至十一年（1638）完成，书中徐光启向国人介绍了欧洲流行于十六世纪第谷（Tycho）「地心说」，并以此为理论体系第一次向国人介绍了地圆的知识理念，并对地球经纬度的概念进行了说明，动摇了长期以来占据中国人意识形态「天圆地方」的观念，其实对于宇宙的认识，哥白尼已经在《天体运行论》（1543年出版）一书中提出了「日心说」，推翻了「地心说」，但书中对十五世纪已经出现的哥白尼（Copernicus)「日心说」天体观没有作介绍，这是《崇祯历书》局限所在。

尽管徐光启在《崇祯历书》中对西方的天道观作了较为系统的介绍，但其理论体系仍未能被民众所普遍接受，甚至对天文历法颇有研究的王夫之对利玛窦为代表的西方天文观表现出来的不屑就很能代表当时普通民众的思想，王夫之《思问录外篇》认为利玛窦「身处大地之中，目力亦与人同，乃倚一远镜之技，死算大地为九万里，使中国有人焉如子瞻、元泽者，曾不足以当其一笑，而百年以来，无有能窥其狂俟者，可叹也！」又论「西洋历家既能测知七曜远近之实，而又窃张子左旋之


% 原书第 138 页（PDF 第 151 页）
\BookPage{138}{151}
\noindent 说以相杂立论，盖西夷之可取者，唯远近测法一术，其他则皆剽袭中国之余绪，而无通理之可守也」1。

由此，我们可以对明代的天道观作作一个评估，明代天文学的成就已经显示，明代负责天文和历法修订的知识分子和上层统治者已经通过西方天文知识体系得知传统的「天圆地方说」是不正确的，对许多自然现象的知识与解释均比民间科学，也即官方所掌握的天道观尽管与十七世纪欧洲相比显得落后，但相比民间仍显进步得多，为什么会出现这种现象呢？

明代一如之前的封建时代，对历法的修订十分重视，一方面是因为历法的修订关乎于指导农业的大任，而农业是封建社会的根本，所以历法不能轻易修改。另一方面是政治上的考虑，认为历法乃关乎天命之大事，是「古法」、「古制」，明代奉行的历法为《大统历》，基本理念是元代的《授时历》，至成化年间，对天气节气等许多现象的预测出现了严重的偏差，以致许多人建议修改历法，但统治者均以「古法未可轻变」、「祖制不可改」12予以否决。所以对于天文的观察，统治者更看重的是「天道」的变化与统治者之间的感应关系，故对于历法的对错民间是没有权利提出质疑的，更不能私自进行修订，据明代的文献我们得知「国初学天文有历禁」，统治者对民间修订历法处罚是「习历者遣戍，造历者殊死」，十分严厉13。

虽然十七世纪上半期由于西方天文学知识的传入，统治者也在一定程度上认识到传统观念的错失，但明代的统治者为了维护「天道」的恒常，不愿将科学知识在民间进行宣传与普及，故造成了民间「天道观」滞后科学成果取得的现象。

十七世纪下半期，历史已经进入清代的治理时代，「天文历法之学在清代风靡一时。上至皇帝，下及布衣，凡知识阶层咸以言天文历法为荣」14。清人的天道观随着西方天文学在中国的传播发生了很大的改变，民间天道观渐多了科学的成分。

\vfill
\begin{BookNotes}
\BookNote{11.}{（清）王夫之撰，《思问录外篇》，载于《船山遗书》，中国古籍出版社，1950年，第42-46页。}
\BookNote{12.}{（清）张廷玉等撰，《明史》，卷31〈历一〉，中华书局，1974年，第519页。}
\BookNote{13.}{（明）沈德符撰，《万历野获编》，卷20〈历法·历学〉，中华书局，1959年，1997年重印，第524页。}
\BookNote{14.}{江晓原，〈17、18世纪中国天文学的三个新特点〉，《自然辩证法通讯》，1988年3期。}
\end{BookNotes}
